\documentclass[11pt,letterpaper]{article}

% ============================================================================
% PACKAGES
% ============================================================================
\usepackage[margin=1in, top=1.5in, headheight=85pt]{geometry}
\usepackage{graphicx}
\usepackage{fancyhdr}
\usepackage{lastpage}
\usepackage{enumitem}
\usepackage{xcolor}
\usepackage{hyperref}
\usepackage{tabularx}

% ============================================================================
% DOCUMENT VARIABLES
% ============================================================================
\newcommand{\UniqueID}{DM-2026-002}
\newcommand{\DocumentDate}{January 21, 2026}
\newcommand{\AuthorName}{SendCUIEmail Development Team}
\newcommand{\AuthorTitle}{Software Engineering}
\newcommand{\ToField}{Project Record}
\newcommand{\SubjectField}{Maximum File Size Limit for Email Attachments}
\newcommand{\OPRField}{SendCUIEmail Project}

% ============================================================================
% DOCUMENT CONFIGURATION
% ============================================================================
\hypersetup{
    colorlinks=false,
    pdfborder={0 0 0},
    pdftitle={DM-2026-002 File Size Limit},
    pdfauthor={SendCUIEmail Project},
}

\setlength{\parindent}{0pt}
\setlength{\parskip}{6pt}

\setlist[enumerate]{leftmargin=0.5in, labelwidth=0.25in, labelsep=0.1in, align=left, itemsep=12pt, topsep=12pt}
\setlist[enumerate,1]{label=\arabic*.}
\setlist[itemize]{leftmargin=0.5in, itemsep=3pt, topsep=6pt}

% ============================================================================
% HEADER AND FOOTER
% ============================================================================
\pagestyle{fancy}
\fancyhf{}
\fancyhead[L]{\textbf{SendCUIEmail}}
\fancyhead[R]{\parbox[b]{3.5in}{\raggedleft\large\textbf{Decision Memorandum}}}
\renewcommand{\headrulewidth}{0pt}
\fancyfoot[L]{\small \UniqueID}
\fancyfoot[C]{\small Page \thepage\ of \pageref{LastPage}}
\fancyfoot[R]{\small \DocumentDate}
\renewcommand{\footrulewidth}{0.4pt}

\newcommand{\dmsection}[1]{\item \textbf{#1}\par}

% ============================================================================
% BEGIN DOCUMENT
% ============================================================================
\begin{document}

\hfill \UniqueID

\vspace{0.25in}

\underline{\hspace{2.5in}}\\[2pt]
\AuthorName\\
\AuthorTitle

\vspace{0.25in}

\textbf{DATE:} \DocumentDate

\textbf{TO:} \ToField

\textbf{SUBJECT:} \SubjectField

\textbf{OFFICE OF PRIMARY RESPONSIBILITY:} \OPRField

\vspace{0.25in}

\begin{enumerate}

% ----------------------------------------------------------------------------
\dmsection{Purpose}

This memorandum establishes the maximum file size limit for SendCUIEmail encrypted attachments and documents the rationale for this limit.

% ----------------------------------------------------------------------------
\dmsection{Decision}

\textbf{Maximum supported attachment size: 10 MB (10,485,760 bytes)}

Files larger than 10 MB should not be sent via SendCUIEmail. For larger files, users should use alternative secure transfer methods.

% ----------------------------------------------------------------------------
\dmsection{Rationale}

\textbf{Email Infrastructure Limitations:}
\begin{itemize}
    \item Most enterprise email systems impose attachment size limits between 10-25 MB
    \item Base64 encoding (used in email transmission) increases file size by approximately 33\%
    \item A 10 MB encrypted file becomes approximately 13.3 MB after Base64 encoding
    \item Many email gateways reject messages larger than 10-15 MB total
\end{itemize}

\textbf{Practical Considerations:}
\begin{itemize}
    \item Email is not designed for large file transfer
    \item Large attachments increase risk of delivery failure
    \item Recipients may have mailbox quotas that prevent receipt
    \item Email filters may quarantine or reject large attachments
\end{itemize}

\textbf{Security Considerations:}
\begin{itemize}
    \item Larger files require more time to encrypt/decrypt, increasing exposure risk
    \item Failed transfers may result in incomplete data on intermediate systems
    \item Log files and email archives may retain copies indefinitely
\end{itemize}

% ----------------------------------------------------------------------------
\dmsection{Alternative Methods for Large Files}

For files exceeding 10 MB, use one of these approved methods:

\begin{itemize}
    \item \textbf{Secure File Transfer Services}: Organization-approved SFTP or cloud storage with encryption at rest
    \item \textbf{Encrypted USB Media}: Physical transfer using FIPS 140-2 validated encrypted drives
    \item \textbf{Secure File Sharing Platforms}: DoD SAFE, MilSuite, or similar approved platforms
    \item \textbf{Split Archives}: For files slightly over the limit, consider splitting into multiple encrypted archives (not recommended for CUI due to complexity)
\end{itemize}

% ----------------------------------------------------------------------------
\dmsection{Implementation}

The file size limit is implemented as follows:

\begin{itemize}
    \item \textbf{Warning}: Files approaching 10 MB will display a warning
    \item \textbf{Test Coverage}: The test suite includes a 10 MB file to verify maximum size handling
    \item \textbf{Documentation}: User documentation clearly states the limit
    \item \textbf{No Hard Block}: The tool does not prevent encryption of larger files, but users are warned
\end{itemize}

\textbf{Note}: This is a ``moral limit'' rather than a technical enforcement. Users are trusted to use appropriate methods for large file transfer.

% ----------------------------------------------------------------------------
\dmsection{Test Data}

The following test files are provided in the \texttt{testdata/} directory:

\begin{tabularx}{\textwidth}{|l|r|X|}
\hline
\textbf{File} & \textbf{Size} & \textbf{Purpose} \\
\hline
\texttt{small\_text.txt} & 138 B & Basic text file \\
\hline
\texttt{empty.txt} & 0 B & Edge case: empty file \\
\hline
\texttt{binary\_sample.bin} & 1 KB & Binary with null bytes \\
\hline
\texttt{README.pdf} & 175 KB & Generated from README.md \\
\hline
\texttt{medium\_file.bin} & 1 MB & Medium binary file \\
\hline
\texttt{large\_file\_10mb.bin} & 10 MB & Maximum size test \\
\hline
\end{tabularx}

% ----------------------------------------------------------------------------
\dmsection{References}

\begin{itemize}
    \item RFC 5321 (SMTP): Recommends message size limits
    \item NIST SP 800-171: CUI handling requirements
    \item DoD Instruction 8170.01: Online Information Management
    \item Organization email policy (varies by deployment)
\end{itemize}

\end{enumerate}

\end{document}
